\begin{lemma}
\label{lemma-pro-isomorphism-iso-Rlim}
Let $(A_n)$ and $(B_n)$ be inverse systems of abelian groups.
A morphism of pro-systems $\varphi : (A_n) \to (B_n)$ determines
maps $\lim A_n \to \lim B_n$ and $R^1\lim A_n \to R^1\lim B_n$.
These maps are isomorphisms if $\varphi$ is a pro-isomorphism.
\end{lemma}

\begin{proof}
Please see
Categories, Example \ref{categories-example-pro-morphism-inverse-systems}
for a discussion of morphisms of pro-systems.
The map $\varphi$ is given by $1 \leq m_1 < m_2 < m_3 < \ldots$
and maps $\varphi_n : A_{m_n} \to B_n$ compatible with transition
maps. Set $\varphi', \varphi'' : \prod A_n \to \prod B_n$
equal to
$\varphi'((a_n)) = (\varphi_n(a_{m_n}))$ and
$$
\varphi''((a_n)) =
(\varphi_n(a_{m_n} + f(a_{m_n + 1}) + \ldots + f(a_{m_{n + 1} - 1})))
$$
where each occurence of $f$ denotes a suitable transition map of
the inverse system $(A_n)$. Then the diagram
$$
\xymatrix{
\prod A_n \ar[r]_\delta \ar[d]^{\varphi'} &
\prod A_n \ar[d]^{\varphi''} \\
\prod B_n \ar[r]^\delta & \prod B_n
}
$$
where the horizontal arrows are as in Lemma \ref{lemma-compute-Rlim}
is commutative. In this way we obtain the desired maps. The construction
is functorial in the sense that if we're given an inverse system $(C_n)$
and $1 \leq m'_1 < m'_2 < m'_3 < \ldots$ and maps $\psi_n : B_{m'_n} \to C_n$
compatible with transitition maps, then
$(\psi \circ \varphi)' = \psi' \circ \varphi'$ and
$(\psi \circ \varphi)'' = \psi'' \circ \varphi''$
where the composition $\psi \circ \varphi$ refers to the integers
$1 \leq m_{m'_1} < m_{m'_2} < \ldots$ and the maps
$\psi_n \circ \varphi_{m'_n} : A_{m_{m'_n}} \to C_n$. We will show
that if $B_n = A_n$ and $\varphi_n$ is the transition map, then
the resulting maps are the identity maps. This will both prove that
the construction is independent of the choice of the representative
$(m_n, \varphi_n)$ of $\varphi$ and the final statement of the lemma.

\medskip\noindent
Thus we let $B_n = A_n$ and $\varphi_n$ be the transition map.
Let $(a_n) \in \lim A_n$ be an element. Then $\varphi'((a_n))$
is the element which has in degree $n$ the image of $a_{m_n}$
which is equal to $a_n$. This proves the statement for $\lim A_n$.
Let $\xi \in R^1\lim A_n$ be the class of the element
$(a_n)$ in $\prod A_n$. Consider the element $(b_n)$ with
$$
b_i = a_i + f(a_{i + 1}) + \ldots + f(a_{m_n - 1})
$$
if $m_{n - 1} < i < m_n$ and $0$ if $i = m_n$ for some $n$.
Then $(a'_n) = (a_n) - \delta((b_n))$ defines the same class
in $R^1\lim A_n$ and a computation shows that $a'_i = 0$
unless $i = m_n$ for some $n$. Thus we may and do assume $a_i = 0$
unless $i = m_n$ for some $n$. Note that
$$
(0, \ldots, -f(a_{m_j}), 0, \ldots, 0, a_{m_j}, 0, \ldots)
$$
with nonzero entries in spots $j$ and $m_j$, is the image under $\delta$ of
$$
c_j = (0, \ldots, 0, f(a_{m_j}), f(a_{m_j}), \ldots, a_{m_j}, 0, \ldots)
$$
with nonzero entries in spots $j + 1, \ldots, m_j$. The sum
$c = \sum c_j$ makes sense in $\prod A_n$. Recalling that
$\varphi''((a_n)) = (a_{m_1}, a_{m_2}, \ldots)$ we see that
$$
(a_n) - \varphi''((a_n)) = \delta(c)
$$
and the proof is complete.
\end{proof}